\documentclass[11pt]{article}
\usepackage{amsmath,amssymb,amsthm,booktabs}
\usepackage[margin=2.4cm]{geometry}
\usepackage[hyphens]{url}\sloppy
\newtheorem{theorem}{Theorem}\newtheorem{lemma}[theorem]{Lemma}\newtheorem{proposition}[theorem]{Proposition}
\newtheorem{corollary}[theorem]{Corollary}
\theoremstyle{definition}\newtheorem{definition}[theorem]{Definition}\newtheorem{remark}[theorem]{Remark}
\newcommand{\Z}{\mathbb Z}
\title{Q597 G1a: the role is a type, and the readout follows the block swap\\
\normalsize PROVE 2026-11-22 c2. \textbf{(Role) is reduced}, with a proof, to the block swap of a \emph{locally defined} journey plus three computed statements. In every diverging $B/C$ pair, exactly one of the two journeys is at a class-1 a-step at its first contact (proved from the mutual handshake and the local classification). Call it $L$. The readout set $I$ is the support of the mixed second difference of the readout, so it belongs to the unordered pair (proved). If $L$ satisfies BS, then $I\subseteq\mathrm{side}(L)$, and $I=\emptyset$ exactly when $L$'s re-run takes an a-step. So (Role) is an \emph{output} of BS for $L$, not an extra input. Retreat: a suffix version of Lemma dinv is proved, and the circularity in (MidEnd) is located exactly.}
\author{Clio}\date{22 November 2026 (c2)}
\begin{document}\maketitle

\section*{Provenance}
\begin{itemize}
\item Code: \texttt{proofs/code-1122c2-Q597/} (new; no script shadows a stdlib module). \texttt{rolesdump.py} runs the census loop of \texttt{code-1122-Q597/pairdich.py} unchanged (through the \texttt{code-1121-Q597/diverge.py} preamble) and pickles every diverging ordered $B/C$ pair and every non-commuting $B/B$ pair with its raw step lists. All other scripts read the pickle offline: \texttt{signtab.py}, \texttt{signdetail.py}, \texttt{orient.py} (Step~0), \texttt{stepside.py} (readout sides, mixed difference), \texttt{perm.py}, \texttt{lformula.py} (BS and the readout formula for $L$), \texttt{share.py} (the shared tile), \texttt{suffix.py} (retreat), and \texttt{examples.py}. Each \texttt{X.py} writes \texttt{X5.out} and \texttt{X6.out}.
\item Sizes: the full census at size 5 (5194 states, \texttt{rolesdump.py 5 4 24}, 4\,min\,10\,s), and the size-6 subsample \texttt{KEEP=3} of 24 chunks (3951 states, 5\,min\,9\,s), as on 11-22.
\item Upstream, all my own files, read first-hand this session: \texttt{2026-11-22-Q597-first-contact.tex} (Lemma~U, Lemma~dinv, Theorem~fc, Theorem~local, Corollary~equiv, Theorem~role, Theorem~a2), \texttt{2026-11-21c2-Q597-sheet.tex} (Theorem~cc, Theorem~red, Lemma~phi), \texttt{2026-11-18-Q597-one-star-SCEL.tex} (Theorem~red and Corollary~side there, the readout formula under BS), and the engine \texttt{code-1112c2-Q585/migrate.py} (\texttt{Board.journey}, the termination test). No external source is used.
\end{itemize}

\section*{Summary}
\begin{enumerate}
\item \textbf{Step 0: is the role a sign?} (\S\ref{sec:step0}). Neither of the brief's potentials gives a sign equal to $I$'s side. $\psi$ has a constant sign over all pairs. $g$ (Lemma~phi) takes $g(R^x_{p})-g(R^y_{p})\in\{0,-1\}$ on every diverging pair, and its value names the journey at a class-1 a-step. That is a \emph{geometric} orientation. It is defined on the empty-$I$ pairs too, so it is not $I$'s side. The step-index sign $\mathrm{sign}(p_x-p_y)$ matched $I$'s side on 521/521 at size 5, but it is $0$ on 3 pairs of the size-6 subsample, one of them one-sided. It is refuted.
\item \textbf{Lemma~\ref{lem:mixed} (proved).} When the boards commute, $X^y-X=Y^x-Y=P_{12}-P_1-P_2+P$. So $I$ is the support of a mixed second difference and is a property of the \emph{unordered} pair. The brief's suspected broken assumption (``$I$ depends on the order of application'') is false. Data: 779/779 and 854/854. Control: it fails on all $B/B$ pairs, 440/440 and 463/463, whose boards do not commute.
\item \textbf{The mutual handshake} (Lemma~\ref{lem:hs}). Each contact step reads the other journey's contact step: $c_0(x)=p_y$ and $c_0(y)=p_x$. The two inequalities $c_0(x)\ge p_y$, $c_0(y)\ge p_x$ are proved, and so is the equivalence of the two equalities with (Share): the two contact hexagons share a tile of $M_0$. (Share) is computed: the two contact hexagons share exactly one tile, a square of $M_0$, in 1219/1219 and 1317/1317 pairs.
\item \textbf{Proposition~\ref{prop:orient} (proved from (Share) and Theorem~local).} In every diverging $B/C$ pair exactly one role, $L$, is at a class-1 a-step $(1,a_T)$ at its contact. $L$'s re-run takes $b$ or $a_{T'}$ (Theorem~local).
\item \textbf{Theorem~\ref{thm:main}: (Role) $\Leftarrow$ (Share) $\wedge$ BS$(L)$ $\wedge$ (Ca) $\wedge$ (N)}, proved. If $L$'s re-run takes $b$, then $I\neq\emptyset$ lies on $L$'s side and $L$ has the (CC) contact. If it takes $a_{T'}$, then $I=\emptyset$. Hence the role on $I$'s side, whenever $I$ is one-sided, is $L$, and it is good. All hypotheses are computed at both sizes: BS$(L)$ 779/779 and 854/854, \emph{including the empty-$I$ pairs}. The lesion (the other role $E$) has BS on 0 of the 521 and 0 of the 506 one-sided pairs.
\item \textbf{Re-indexing} (Remark~\ref{rem:reindex}). (CC) holds for any $B/C$ role whose contact is $(1,a_T)/(1,b)$, by Theorem~local and Corollary~equiv alone. So the G1 program can be run on $L$, which is defined locally, instead of on ``the role on $I$'s side''. (Role) then comes out of BS$(L)$ through Theorem~\ref{thm:main}. Two things remain for this: a line-by-line check that Theorem~cc's chain never uses $I$, and the two-step swap (case~a) as a separate statement.
\item \textbf{Retreat: (MidEnd)} (\S\ref{sec:midend}). Lemma~\ref{lem:suffix} (suffix dinv) is proved: if the journey and its re-run stand on the same rhombus, the re-run copies the journey for as long as the journey's hexagons avoid the board difference. It needs no hypothesis on the re-run's side, because of Lemma~U. So (MidEnd) $\Leftarrow$ (Sh3) $\wedge$ (Resync) $\wedge$ (Stale) $\wedge$ (T). The circularity is exactly (Resync): the re-run stands on the journey's rhombus at step $p+m+1$. That is the rejoin clause of BS.
\end{enumerate}

\section{Setting}\label{sec:set}
As in the first-contact tex \S1. For an ordered pair of corners $(x,y)$ of nest~$A$ in a census state $M_0$, $J_x$ is $x$'s journey toward $B$ and $J_y$ is $y$'s journey toward $C$, both on $M_0$. The \emph{role} $x$ is the triple ($A=J_x$, its re-run $A'$ on $M_0\cdot J_y$, $J_o=J_y$), and symmetrically for the role $y$. $p$ is the length of the longest common prefix of $A$ and $A'$, $W=M_0\setminus M_0\cdot J_o$, and $c_0$ is the least $J_o$ step that removes a tile of $H_p\cap W$. A pair \emph{diverges} if some role has a middle. Readouts: $P=\mathrm{pq}(M_0)$, $P_1=\mathrm{pq}(M_0\cdot J_x)$, $P_2=\mathrm{pq}(M_0\cdot J_y)$, $P_{12}$ and $P_{21}$ after both. $X=P_1-P$, $Y=P_2-P$, $Y^x=P_{12}-P_1$, $X^y=P_{21}-P_2$, and
\[I=\{k:\ Y^x(k)\ne Y(k)\ \text{or}\ X^y(k)\ne X(k)\}.\]
Each key $k$ has a side $k_0\in\{0,1\}$. $I$ is \emph{one-sided on side $s$} if $\emptyset\ne I\subseteq\{k_0=s\}$. The clean journey of a clean record is the role on $I$'s side: $x$ for side~1, $y$ for side~0. Types are $(\text{class},\text{type})$ as in \texttt{diverge.py}: $a1$ means the rhombus displacement $\alpha$ has zero $\Lambda_0$-part, and $a0$ means zero $\Lambda_1$-part. Write $a_B=a1$, $a_C=a0$.

\section{Step 0: is the role a sign?}\label{sec:step0}
Preconditions, printed first (\texttt{signtab5.out}, \texttt{signtab6.out}). 779 diverging ordered $B/C$ pairs at size 5 (of 1558; the other 779 diverge on neither side), and in every diverging pair \emph{both} roles diverge. $\mathrm{fc}=p$ on both roles. Sides 268 (1), 253 (0), 258 (empty), none two-sided. $B/B$: 440 non-commuting pairs, both roles diverging. Size-6 subsample: 854 diverging $B/C$ (210, 296, 348) and 463 $B/B$.

\begin{center}\small\begin{tabular}{lrrrr}\toprule
size 5 (size 6) & side 1 & side 0 & $I=\emptyset$ & $B/B$\\\midrule
$\Delta g=g(R^x_{p_x})-g(R^y_{p_y})=0$ & 268 (210) & 0 (0) & 211 (295) & 440 (463)\\
$\Delta g=-1$ & 0 (0) & 253 (296) & 47 (53) & 0 (0)\\
$\mathrm{sign}\,\Delta\psi_0$, $\mathrm{sign}\,\Delta\psi_1$ & $+,-$ & $+,-$ & $+,-$ & $+,-$\\
$\mathrm{sign}(p_x-p_y)=+1$ & 268 (210) & 0 (0) & 211 (294) & 440 (463)\\
$\mathrm{sign}(p_x-p_y)=0$ & 0 (0) & 0 (\textbf{1}) & 0 (2) & 0 (0)\\
$\mathrm{sign}(p_x-p_y)=-1$ & 0 (0) & 253 (295) & 47 (52) & 0 (0)\\
\bottomrule\end{tabular}\end{center}
Here $\psi$ is the no-repeat potential with $T$'s heights, $(\min h_0,\min h_1)$ over the contact rhombus. Both orientations are in the table: with $y\to C$ first, every sign flips.
\begin{itemize}
\item $\psi$ has a constant sign: it carries no information about the pair.
\item $\Delta g$ takes two values and never $+1$. It is not $I$'s side, because both values occur on empty-$I$ pairs. What it does name is the journey at a class-1 a-step (Proposition~\ref{prop:orient}): $\Delta g=0$ iff that journey is $x$, on all 779 and 854 pairs. This is a property of the local contact datum, since the offset of the two contact rhombi is part of it (Theorem~local).
\item The step-index sign matched $I$'s side on 521/521 at size 5. On the size-6 subsample it is $0$ on one side-0 pair and two empty pairs. \textbf{Refuted.} The ``later contact'' reading is a size-5 coincidence.
\item \textbf{Role-swap lesion} (\texttt{orient.py}): calling the other role $L$ makes the prediction of Theorem~\ref{thm:main} fail on \emph{all} one-sided pairs (521/521 and 506/506). It agrees only on 9 (20) empty-$I$ pairs, where the swapped role's re-run happens to be an a-step. $B/B$ control: the class-1 a-step role is BBy in 440/440 and 463/463, and it is the role with $(\mathrm{Sheet}_p)\wedge|H_p\cap W|=2$. Swapping gives ``not good'' on all of them. The instrument reads the roles.
\end{itemize}
So the answer to Step~0 is: \emph{the role is not a sign of a potential, it is a type.} One journey of each crossing stands at a class-1 a-step when it meets the other. Which side $I$ falls on is then decided by that journey's re-run (\S\ref{sec:main}).

\section{The pair, symmetrically}\label{sec:pair}
\begin{lemma}[the readout set is a mixed difference]\label{lem:mixed}
Suppose $M_0\cdot J_x\cdot J_y^{\,x}=M_0\cdot J_y\cdot J_x^{\,y}$ as boards (the re-runs are the journeys on the moved boards). Then $X^y-X=Y^x-Y=P_{12}-P_1-P_2+P=:\Delta$, and $I=\operatorname{supp}\Delta$. In particular $I$, and hence its side, does not depend on which journey is applied first.
\end{lemma}
\begin{proof}
The readout is a function of the board, so $P_{21}=P_{12}$. Then $X^y-X=(P_{21}-P_2)-(P_1-P)=P_{12}-P_1-P_2+P$, and $Y^x-Y=(P_{12}-P_1)-(P_2-P)$ is the same expression. $k\in I$ iff one of the two differences is nonzero at $k$, and they are equal. The differences are taken keywise. That is legitimate because every step preserves the key set of the readout (engine check: 37\,268 steps at size 5, 43\,694 on the subsample, every step changes the values of exactly two keys and no key appears or disappears; \texttt{stepside.py}).
\end{proof}
\noindent Data: the boards commute in all 1558 ordered $B/C$ pairs (pairdich), and $X^y-X=Y^x-Y$ with support $I$ in 779/779 and 854/854 diverging pairs. \emph{Control:} on the non-commuting $B/B$ pairs the identity fails in 440/440 and 463/463, as it must.

\begin{lemma}[the handshake is mutual, half]\label{lem:hs}
Let $J_1$, $J_2$ be terminating journeys on $M_0$, let $W_i$ be the tiles of $M_0$ removed by $J_i$, let $p_i$ be the first step of $J_i$ whose hexagon meets $W_{3-i}$, and let $c_0(i)$ be the least step of $J_{3-i}$ that removes a tile of $H^i_{p_i}\cap W_{3-i}$. Then:
\begin{enumerate}
\item $c_0(1)\ge p_2$ and $c_0(2)\ge p_1$;
\item $c_0(1)=p_2$ $\iff$ $c_0(2)=p_1$ $\iff$ \textup{(Share)}: $H^1_{p_1}\cap H^2_{p_2}\cap M_0\ne\emptyset$.
\end{enumerate}
\end{lemma}
\begin{proof}
By Lemma~nr (no re-creation) and no revisit, $W_i=M_0\setminus M_0\cdot J_i$, and each tile of $W_i$ is removed by $J_i$ exactly once. Write $\mathrm{rem}_i(u)$ for that step.
(1) Let $u\in H^1_{p_1}\cap W_2$ with $\mathrm{rem}_2(u)=c_0(1)$. Then $u\in M_0$ and $J_1$ removes $u$ at $p_1$, so $u\in W_1$. So $J_2$'s step $c_0(1)$ removes a tile of $W_1$, and its hexagon meets $W_1$. By minimality of $p_2$, $p_2\le c_0(1)$. The other inequality is the same with the indices exchanged.
(2) If $c_0(2)=p_1$, take $v\in H^2_{p_2}\cap W_1$ with $\mathrm{rem}_1(v)=p_1$. Then $v\in H^1_{p_1}\cap M_0$, and $J_2$ removes $v$ at $p_2$, so $v\in W_2$ and $c_0(1)\le\mathrm{rem}_2(v)=p_2$. With (1), $c_0(1)=p_2$. The converse is symmetric. If $c_0(1)=p_2$, the tile $u$ of (1) lies in $H^1_{p_1}\cap H^2_{p_2}\cap M_0$. Conversely, a tile $w$ there has $w\in W_1\cap W_2$ with $\mathrm{rem}_1(w)=p_1$ and $\mathrm{rem}_2(w)=p_2$, so $c_0(1)\le p_2$, and (1) gives equality.
\end{proof}
\noindent\textbf{Computed (Share)} (\texttt{share.py}): in every diverging pair the two contact hexagons share \emph{exactly one} tile, and it is a square of $M_0$. No non-$M_0$ tile is shared. 779 + 440 at size 5, 854 + 463 on the subsample, over all eight $B/C$ type combinations and $B/B$. The mutual equalities $c_0(x)=p_y$, $c_0(y)=p_x$ hold in all of these pairs (\texttt{signtab.py}, \texttt{MUTUAL}). (Share) is the 11-21 c2 handshake ``$H_p\cap H^o_{c_0}=\{\text{square}\}$'' seen from both sides at once.

\begin{proposition}[orientation]\label{prop:orient}
Let $(x,y)$ be a diverging $B/C$ pair in which both roles diverge and satisfy \textup{(FC$\le$)}, and assume \textup{(Share)}. Then exactly one role is at a class-1 a-step of its own target at its contact: either $x$ at $(1,a1)$ or $y$ at $(1,a0)$. Call it $L$ and the other role $E$. $L$'s re-run takes $(1,b)$ or $(1,a_{T'})$.
\end{proposition}
\begin{proof}
By Theorem~local of the first-contact tex, the contact datum of each role is a survivor of \texttt{ccloc.py}. For the role $x$ (target $B$, $T'=C$) the survivors have $A$-type and first-read $J_o$-type
\[(1,a1):\ (0,a0)\text{ or }(1,b);\qquad (0,a1):\ (1,a0);\qquad (1,b):\ (1,a0).\]
The first-read $J_o$ step is $c_0(x)$, and by (Share) and Lemma~\ref{lem:hs}, $c_0(x)=p_y$. So the first-read type is the type of $y$'s step at its own contact. Hence, if $x$ is at $(1,a1)$, then $y$ is at $(0,a0)$ or $(1,b)$, not at $(1,a0)$. If $x$ is not at $(1,a1)$, then $y$ is at $(1,a0)$. Exactly one of ``$x$ at $(1,a1)$'' and ``$y$ at $(1,a0)$'' holds. The re-run types are read off the same table: with $A$ at $(1,a1)$ the re-run is $(1,b)$ or $(1,a0)$. The $\Sigma$-image gives the table for $y$, with $(1,a0)$ re-run $(1,b)$ or $(1,a1)$, and the same argument read from $y$'s side is consistent with this.
\end{proof}
\noindent\emph{Data:} exactly one class-1 a-step role in 779/779 and 854/854 $B/C$ pairs, and 440/440, 463/463 $B/B$ (\texttt{orient.py}, \texttt{PRE}). The eight realised $(L,E)$ type combinations are listed in \texttt{share5.out}.

\section{(Role) from the block swap of $L$}\label{sec:main}
For a role with middle $m_A=a\cdot C$ and re-run middle $m_B$, \textbf{BS} is the 11-18 block swap: $m_B=(C-\alpha)\cdot(a+\gamma)$ as rigid $\Z^4$ translates, with $\alpha$ the rhombus displacement of $a$ and $\gamma$ that of $C$ (from the rhombus before $C$'s first step to the rhombus after its last).

\begin{lemma}[readout sides of the step types]\label{lem:side}
An $a0$-step changes the readout only on side-0 keys, an $a1$-step only on side-1 keys, and a $b$-step on one key of each side.
\end{lemma}
\noindent\emph{Status: engine check}, not a proof by hand. Over the 37\,268 and 43\,694 steps of \texttt{stepside.py} (all $A$ and re-run steps of the dumped pairs, both targets, both classes), with no exception. It is a statement about the twelve good configurations (Lemma~word), so a finite check on the configuration table would make it exhaustive. I did not run that check.

\begin{theorem}[(Role) from BS of $L$]\label{thm:main}
Let $(x,y)$ be a diverging $B/C$ pair satisfying \textup{(FC$\le$)} on both roles and \textup{(Share)}, and let $L$ be as in Proposition~\ref{prop:orient}, with target $T$, $s=1$ if $T=B$ and $s=0$ if $T=C$. Assume:
\begin{itemize}
\item[\textbf{BS$(L)$}] $L$'s middle and re-run middle satisfy BS;
\item[\textbf{(Ca)}] if $L$'s re-run takes $(1,a_{T'})$, every step of $C$ is an $a_{T'}$-step;
\item[\textbf{(N)}] if $L$'s re-run takes $(1,b)$, then $\Delta\ne0$.
\end{itemize}
Then $I=\emptyset$ when $L$'s re-run takes $a_{T'}$, and $I$ is one-sided on side $s$ when it takes $b$. Consequently, \textup{(Role)} holds: if $I$ is one-sided, the role on $I$'s side is $L$. Its contact is $(1,a_T)/(1,b)$, and so it has $(\mathrm{Sheet}_p)$, $|H_p\cap W|=2$ and \textup{(CC)}.
\end{theorem}
\begin{proof}
Let $\delta_\sigma$ be the readout change of a step $\sigma$, and $T_v$ the shift of keys by $v$ (side-$i$ keys by the $\Lambda_i$-block of $v$). By Lemma~\ref{lem:mixed} we may compute $\Delta$ in the role $L$: $\Delta=\sum_{m_B}\delta-\sum_{m_A}\delta$, since the common prefix and suffix steps are identical and cancel. Under BS this is Theorem~red of the 11-18 SCEL tex:
\[\Delta=\sum_{\sigma\in C}\bigl(T_{-\alpha}\delta_\sigma-\delta_\sigma\bigr)+\bigl(T_\gamma\delta_a-\delta_a\bigr).\]
Here $a$ is $L$'s contact step, of type $a_T$. By the definition of the type, $\alpha\in\Lambda_s$: its $\Lambda_{1-s}$-block is $0$. By Lemma~\ref{lem:side}, $\delta_a$ is supported on side $s$.
\emph{Re-run $b$.} $T_{-\alpha}$ fixes side-$(1-s)$ keys, so each $T_{-\alpha}\delta_\sigma-\delta_\sigma$ vanishes on side $1-s$. $T_\gamma\delta_a-\delta_a$ is supported on side $s$ because $\delta_a$ is. So $\Delta$ is supported on side $s$, and by (N) it is nonzero. By Lemma~\ref{lem:mixed}, $I$ is one-sided on side $s$.
\emph{Re-run $a_{T'}$.} By BS the re-run's first middle step is $C[0]-\alpha$, of the same type as $C[0]$. By (Ca) every $\sigma\in C$ is an $a_{T'}$-step, so $\delta_\sigma$ lives on side $1-s$, which $T_{-\alpha}$ fixes. So the first sum vanishes. $\gamma$ is the sum of the displacements of the steps of $C$, each in $\Lambda_{1-s}$, so $T_\gamma$ fixes side-$s$ keys, and $T_\gamma\delta_a=\delta_a$. Hence $\Delta=0$ and $I=\emptyset$.
\emph{(Role).} If $I$ is one-sided on side $s'$, then $I\ne\emptyset$. So $L$'s re-run is not $a_{T'}$, and by Proposition~\ref{prop:orient} it is $b$. Then $s'=s$, i.e.\ the role on $I$'s side ($x$ for side~1, $y$ for side~0) is $L$. Its contact is $(1,a_T)$ with re-run $(1,b)$. In the table of Theorem~local the rows with these types are exactly the two (CC) rows, with $(\mathrm{Sheet}_p)$ and $|H_p\cap W|=2$ (Corollary~equiv).
\end{proof}

\paragraph{Computed hypotheses and controls} (\texttt{lformula.py}, \texttt{perm.py}).
\begin{center}\small\begin{tabular}{lrrrr}\toprule
size 5 (size 6) & side 1 & side 0 & $I=\emptyset$, $L=x$ & $I=\emptyset$, $L=y$\\\midrule
$L$'s re-run & $b$: 268 (210) & $b$: 253 (296) & $a0$: 211 (295) & $a1$: 47 (53)\\
BS$(L)$ & 268 (210) & 253 (296) & 211 (295) & 47 (53)\\
formula of Thm red $=\Delta$ & 268 (210) & 253 (296) & 211 (295) & 47 (53)\\
$\alpha\in\Lambda_s$, $\delta_a$ on side $s$ & all & all & all & all\\
(Ca): $C$ & --- & --- & one $(1,a0)$ step & one $(1,a1)$ step\\
(N): $\Delta\ne0$, on side $s$ & 268 (210) & 253 (296) & $\Delta=0$ & $\Delta=0$\\
\midrule
\emph{lesion:} BS$(E)$ & 0 (0) & 0 (0) & 204 of 258 (244 of 348) & \\
\bottomrule\end{tabular}\end{center}
So in the empty-$I$ pairs $L$'s middle is the two-step word $(1,a_T)(1,a_{T'})$ and the re-run swaps it. These are the commuting a-steps on the two sides (\texttt{perm5.out}: \texttt{1a1 1a0 -> 1a0 1a1} 211 times, \texttt{1a0 1a1 -> 1a1 1a0} 47 times). On one-sided pairs the other role $E$ never has BS, as the 11-19 file found for ``the other journey''. On empty-$I$ pairs it often does, and there its formula also gives $\Delta=0$. On $B/B$, $L$ (BBy) fails BS in 440/440 and 463/463, and $E$ in all. That is the control: Theorem~\ref{thm:main} is a valid implication, and on $B/B$ its hypothesis fails.

\begin{remark}[re-indexing the program by $L$]\label{rem:reindex}
Theorem~role of the first-contact tex used (Role) for exactly one thing: to get (CC) for the clean journey. For a $B/C$ role whose contact is $(1,a_T)/(1,b)$, (CC) holds by Theorem~local and Corollary~equiv alone, because those are the two (CC) rows. So define a \emph{geometric record} as a diverging $B/C$ pair together with its role $L$, when $L$'s re-run takes $b$. Then:
\begin{itemize}
\item (FC$\le$) holds (Theorem~fc), and (CC) holds by definition and the finite check;
\item Theorem~cc of the sheet tex gives (A1) $\wedge$ (Sheet) for $L$, \emph{provided} its proof uses ``clean'' only through (FC$\le$) $\wedge$ (CC). I re-read Theorem~cc, Theorem~red, Lemma~bc and the shadow induction in the sheet tex, and none of them mentions $I$. The c1 Theorem~red ((iii)$\Rightarrow$(i)) and Proposition~hs were not re-read this session. \textbf{This check is owed.}
\item If the remaining program (G1a-A $\to$ G1b $\to$ G1c $\to$ WS $\to$ CT $\to$ BS) is likewise $I$-free, it proves BS$(L)$ for geometric records. Then Theorem~\ref{thm:main} returns (Role) as a corollary, instead of taking it as a hypothesis.
\end{itemize}
What the re-indexed program does \emph{not} cover is case~a: $L$'s re-run takes $a_{T'}$, the empty-$I$ pairs. There Theorem~\ref{thm:main} needs BS$(L)$ for a two-step middle. That is a separate, smaller statement (\S\ref{sec:gaps}).
\end{remark}

\section{Retreat: where (MidEnd) is circular}\label{sec:midend}
\begin{lemma}[suffix dinv]\label{lem:suffix}
Let $A$ be a terminating journey on boards $X_j$ and $A'$ a terminating journey on boards $Y_j$, toward the same target, in Purbhoo mosaics (so Lemma~U applies to both). Suppose that at some step $k$ both stand on the same rhombus $R$. Put $D=X_k\setminus Y_k$, and assume:
\begin{itemize}
\item[(i)] no hexagon of $A$ at a step $j\ge k$ meets $D$;
\item[(T)] at every rhombus of $A$ from step $k$ on, the termination test gives the same answer for $A$ and $A'$.
\end{itemize}
Then $A'$ agrees with $A$ from step $k$ on, step for step, and both stop at the same step.
\end{lemma}
\begin{proof}
Induction on $j\ge k$. Suppose $A'[k{:}j]=A[k{:}j]$, so $A'$ stands on $R_j$. The difference $X_j\setminus Y_j$ is still $D$, by the argument of Lemma~dinv: an identical step removes the same tiles from both boards and creates the same tiles in the region they covered. That argument needs the removed hexagon to be on both boards, which the induction provides.
If $A$ stops at $R_j$ (the termination test fires), then by (T) so does $A'$.
Otherwise $A$ takes $H_j\subseteq X_j$. By (i), $H_j\cap D=\emptyset$, so $H_j\subseteq Y_j$. By (T), $A'$ does not stop, and since it terminates it is not stuck. So by Lemma~U, $A'$ takes the unique good configuration at $R_j$ present on $Y_j$. $H_j$ is a good configuration at $R_j$ (Proposition~good of the no-repeat file) and is present on $Y_j$, so $A'$ takes $H_j$. The rotation is fixed by $H_j$, $R_j$ and the target (S3), so the steps are identical.
\end{proof}
\noindent No hypothesis is needed on $Y_k\setminus X_k$: a created tile on the re-run's side cannot divert it, because a second present configuration would stick the engine. This is determinism again, as in Theorem~fc.
(T) is a real hypothesis. \texttt{Board.journey} stops when the rhombus fits into another nest as an ideal extension, and the two boards' nests differ by $J_o$'s added cell.

\paragraph{(MidEnd) through the lemma.} Take $k=p+m+1$, the step after the last middle step that (MidEnd) names. Then (MidEnd) is ``$A$ and $A'$ agree from $k$ on, and differ at $p+m$''. The difference $D=X_k\setminus Y_k$ consists of three kinds of tile:
\begin{enumerate}
\item[(W)] wake tiles not yet removed by $A$;
\item[(C)] tiles created by $A$'s middle and not by $A'$'s;
\item[(M)] tiles of $M_0\cap Mx$ removed by $A'$'s middle but not by $A$'s.
\end{enumerate}
Proposition~sh3 of the sheet tex (from (CC); proved, past half by finite check) says $A$ reads no wake tile after $p+m$. That disposes of (W). So
\[\text{(MidEnd)}\ \Leftarrow\ \text{(Sh3)}\ \wedge\ \textbf{(Resync)}\ R'_k=R_k\ \wedge\ \textbf{(Stale)}\ \text{$A$ never reads a (C) or (M) tile after $k$}\ \wedge\ \text{(T)}\]
(together with $A[p+m]\ne A'[p+m]$, the easy half). Data (\texttt{suffix.py}): at $k=p+|m_A|$ on all $L$-records with re-run $b$, $R_k=R'_k$ in 341 (329), and both journeys end at $k$ in 180 (177). Neither journey's later hexagons meet its board difference, in 521/521 and 506/506. $D\not\subseteq W$ in every record, so (Stale) has content. \emph{Lesion:} one step earlier ($k-1$), the rhombi differ and the hexagons meet $D$ in 521/521 and 506/506.

\textbf{The circularity, stated exactly.} (Resync) is the rejoin clause of BS: $m_B$ ends where $m_A$ ends. It is CT2 / WS$_{\rm last}$ of the 11-19 file. So Theorem~a2 of the first-contact tex, (CC) $\wedge$ (MidEnd) $\Rightarrow$ (A2), uses the rejoin. If WS is proved from (A2), the program is circular at (Resync) and nowhere else. (Sh3) and Lemma~\ref{lem:suffix} are not circular.
Why the circle cannot simply be cut with (Sh3): in Theorem~a2, (MidEnd) is used to bound the step $k$ of an (A)-tile by $p+m$, after which the shadow puts the tile in $W$. Without (MidEnd), the middle might extend past $p+m$. An (A)-tile $u$ at such a step lies in $A$'s hexagon (Lemma~aread), so by (Sh3) it is not in $W$. It is then a tile of $M_0\cap Mx$ or a prefix-made tile, and the parity argument says nothing about it. ``The middle is empty after $p+m$'' is (Resync) plus Lemma~\ref{lem:suffix}. So (A2) $\Leftarrow$ (CC) $\wedge$ (Resync) $\wedge$ (Stale) $\wedge$ (T), and I see no route to (A2) that avoids the rejoin.

\section{Where the proof stops}\label{sec:gaps}
\begin{description}
\item[Proved by hand.] Lemma~\ref{lem:mixed}. Lemma~\ref{lem:hs} (both inequalities, and the equivalence with (Share)). Proposition~\ref{prop:orient} from (Share) and Theorem~local. Theorem~\ref{thm:main} from its three hypotheses, Lemma~\ref{lem:side} and the 11-18 Theorem~red. Lemma~\ref{lem:suffix} from Lemma~U, Lemma~dinv's argument and (T).
\item[Engine checks / finite checks used.] Theorem~local and Corollary~equiv (exhaustive finite checks, 11-22). Lemma~\ref{lem:side} (engine check, 81\,k steps; a finite check on the configuration table is owed). Key preservation in Lemma~\ref{lem:mixed} (same engine check).
\item[Computed only.] (Share), 1219/1219 and 1317/1317. BS$(L)$ in all diverging $B/C$ pairs (779/779, 854/854), including (Ca) and (N). (Resync), (Stale), (T) at the rejoin (521/521, 506/506).
\item[Refuted.] The step-index sign $\mathrm{sign}(p_x-p_y)$ as $I$'s side (one side-0 pair with $p_x=p_y$ at size 6). ``$I$ depends on the order of application'' (Lemma~\ref{lem:mixed}).
\item[Open.] (Share): why the two contact hexagons share their square. Case~a of BS$(L)$: the two-step swap $(1,a_T)(1,a_{T'})\to(1,a_{T'})(1,a_T)$ with exact translation. (N): nonvanishing of the readout formula in case~b. The $I$-freeness of the upstream chain (Remark~\ref{rem:reindex}). (Resync) for (MidEnd).
\item[Hypotheses.] Termination, Purbhoo mosaic, the engine rule (S1)--(S4), as before.
\end{description}

\section*{Status}
\begin{itemize}
\item \textbf{Proved:} Lemma~\ref{lem:mixed} ($I$ is a mixed second difference, order-free). Lemma~\ref{lem:hs} (half of the mutual handshake; equivalence with (Share)). Proposition~\ref{prop:orient} (orientation, given (Share)). Theorem~\ref{thm:main}: (Share) $\wedge$ BS$(L)$ $\wedge$ (Ca) $\wedge$ (N) $\Rightarrow$ (Role). Lemma~\ref{lem:suffix} (suffix dinv), and the exact location of the (MidEnd) circularity: (Resync).
\item \textbf{Refuted:} the step-index sign; the order-dependence of $I$.
\item \textbf{Computed:} (Share); BS$(L)$, (Ca), (N) at both sizes; Step~0 tables; the suffix hypotheses at the rejoin.
\item \textbf{Not proved:} (Role) unconditionally; (Share); BS$(L)$; (MidEnd).
\end{itemize}
\end{document}
